\documentclass[a4paper,12pt]{article}
\usepackage{assignmentCian}
\author{Cian Booth: \textbf{s1256705}}
\graphicspath{{img/},{graphs/}}
%\subtitle{Solving Where's Wally Puzzles Computationally}
\def\biblio {
  \bibliographystyle{ieeetr}
  \bibliography{/home/cian/Project/HPCWheresWallySolver/writeup/whereswally}
}

\begin{document}
  \def\biblio{}
  \begin{minipage}{0.75\textwidth}
    \resizebox{72pt}{!} {
      \includegraphics{logos/logo_black}
    }
  \end{minipage}
  \begin{minipage}{0.2\textwidth}
    \resizebox{72pt}{!} {
      \includegraphics{logos/eucrest}
    }
  \end{minipage}
  \begin{center}
    \textbf{\Huge Finding Wally}\\
    {\Large Parallel Object Recognition}\\
    \vfill
    \includegraphics[width=0.4\textwidth]{wallylogo}
    \vfill
    {\huge \textbf{Cian Booth}}\\
    {\large MSc. Dissertation Report}\\
    \vspace{1em}
  \end{center}
  \pagestyle{empty}
  \newpage
    \mbox{}
  \newpage
  \abstract{
    A report on the potential of parallel computer vision for everyday use.
    Discussed within are algorithms and libraries that could make this a possibility.
    Special focus is put on the computer vision library OpenCV.
    Where's Wally? puzzles are used as a test case for parallel object recognition.
  }
  \pagebreak
  \tableofcontents
  \listoffigures
  \listoftables
  \listoflistings
  \pagestyle{empty}
  \pagebreak
  \section*{Acknowledgements}
    Many thanks go to my supervisor, Alistair Grant, for directing me to such an interesting project, and keeping me on the rails.
  \pagestyle{empty}
  \pagebreak
  \pagestyle{plain}
  \setcounter{page}{1}
  \subfile{introduction}
  \subfile{background_information}
  \subfile{algorithms}
  \subfile{implementation}
  \subfile{results}
  \subfile{conclusion}
  \pagebreak
  \pagestyle{empty}
  \nocite{opencvcheatsheet}
  \bibliographystyle{ieeetr}
  \bibliography{whereswally}
\end{document}
